%
% -------------------- LaTeX-Template für Abschlussarbeiten --------------------
%
% Autor:    Simon Lehmann <simon.lehmann@gmx.de>
% Version:  2.0 (2009-09-10)
% 
% Versionshistorie: 
% | 2.0 | 2009-09-10 | Umbenennung FH Wiesbaden in Hochschule RheinMain
% | 1.1 | 2008-11-15 | Layout verbessert
% | 1.0 | 2008-06-19 | Urversion
% 
% Das Template verwendet die KOMA-Script Dokumentenklassen und setzt einige
% LaTeX Pakete voraus, die unter Umständen nicht im Standardumfang der lokalen
% LaTeX Distribution enthalten sind. Erhältlich sind diese Pakete über das CTAN
% (http://www.ctan.org/). Zur Installation siehe die Dokumentation des jeweili-
% gen Pakets und der LaTeX Distribution.
%
% Weiterhin ist das Template auf die Verwendung von UTF-8 als Zeichenkodierung
% ausgelegt. Dadurch können Umlaute und andere Sonderzeichen direkt verwendet
% werden.
%
% Es sind für die variablen Felder auf der Titelseite und den formalen Erklär-
% ungen Befehle definiert, mit denen diese gesetzt werden können. Siehe dazu
% weiter unten den Abschnitt 'Vorspann'.
% 
% Für die Glossarerstellung wird das glossaries Paket als Nachfolger von
% glossary verwendet. Das Paket ist leider noch nicht in den meisten Distribu-
% tionen von Haus aus verfügbar und muss eventuell manuell über das CTAN (s.o.)
% bezogen werden. Doku siehe <http://tug.ctan.org/tex-archive/macros/latex/
% contrib/glossaries/glossaries-manual.html>
%
%
% Zum Übersetzen sind in der Regel folgende Durchläufe nötig:
%  1. latex arbeit
%  2. bibtex arbeit (nur wenn schon Zitate im Text vorhanden sind)
%  3. makeglossaries arbeit
%  4. latex arbeit
%  5. latex arbeit (damit Referenzen und Verzeichnisse stimmen)
%
% --------------------------- Dokumenteneinrichtung ----------------------------
% Einsteigern wird empfohlen direkt zum Abschnitt 'Vorspann' weiter unten zu
% springen und die dokumentenweiten Einstellungen zunächst nicht zu verändern.
%
\documentclass[ngerman,a4paper,11pt,twoside,openright,cleardoubleempty,halfparskip]{scrreprt}

% Erforderliche Pakete
\usepackage[utf8]{inputenc}
\usepackage{babel}
\usepackage{graphicx}
\usepackage{fancyhdr}
\usepackage{multirow}
\usepackage{ifthen}
\usepackage{calc}
\usepackage{tabularx}

% Laden von weiteren, nützlichen Paketen (je nach Bedarf)
%\usepackage{subfigure} 
%\usepackage{varioref}
%\usepackage{framed}
%\usepackage{longtable}
%\usepackage{floatflt}
%\usepackage{amsmath}
%\usepackage{url}
% ...

% Pakete für Verzeichnisse. Diese sind hier aufgeführt, da mindestens das
% 'glossaries' Paket nach z.B. hyperref geladen werden muss.
\usepackage{bibgerm}
\usepackage[toc,acronym]{glossaries}

% Glossar(e) laden und erstellen
\makeglossaries
\loadglsentries{glossar} % Name der Glossardatei (ohne .tex)

% Befehlsdefinitionen
%
% Befehlsdefinitionen für das LaTeX Template
%
% Autor:	Simon Lehmann <simon.lehmann@gmx.de>
% Version:	1.0 (2008-06-19)
%

% Sorgt dafür, dass Kapitel immer auf einer rechten Seite beginnen.
\newcommand{\clearemptydoublepage}{\newpage{\pagestyle{empty}\cleardoublepage}}

% Befehle zum Setzen und Abrufen des Ortsnamen, der in der Erklärung verwendet
% wird
\newsavebox{\citybox}
\newcommand{\city}[1]{\sbox{\citybox}{#1}}
\newcommand{\thecity}{\usebox{\citybox}}

% Befehle zum Setzen und Abrufen des Referenten
\newsavebox{\supervisorbox}
\newcommand{\supervisor}[1]{\sbox{\supervisorbox}{#1}}
\newcommand{\thesupervisor}{\usebox{\supervisorbox}}

% Befehle zum Setzen und Abrufen des Korreferenten
\newsavebox{\cosupervisorbox}
\newcommand{\cosupervisor}[1]{\sbox{\cosupervisorbox}{#1}}
\newcommand{\thecosupervisor}{\usebox{\cosupervisorbox}}

\makeatletter % Zugriff auf interne Befehle
% Autorenname, der mit \author{...} gesetzt wurde
\newcommand{\theauthor}{\@author}

% Datum, das mit \date{...} gesetzt wurde
\newcommand{\thedate}{\@date}

% Titel der Arbeit, der mit \title{...} gesetzt wurde
\newcommand{\thetitle}{\@title}
\makeatother

% Flags für Verbreitungsformen
\newboolean{publishlibrary}
\setboolean{publishlibrary}{false}
\newcommand{\publishlibrary}{\setboolean{publishlibrary}{true}}

\newboolean{publishtitle}
\setboolean{publishtitle}{false}
\newcommand{\publishtitle}{\setboolean{publishtitle}{true}}

\newboolean{publishdocument}
\setboolean{publishdocument}{false}
\newcommand{\publishdocument}{\setboolean{publishdocument}{true}}


% Globale Formatierungseinstellungen
\renewcommand{\encodingdefault}{OT1}
\renewcommand{\familydefault}{cmss} % Schriftfamilie auf Sans Serif
\renewcommand{\glsdisplayfirst}[4]{\textit{#1#4}}
\setcapindent{1em}

% Anpassung der Ränder und Breitenverhältnisse
% Der letzte Wert in oddsidemargin bestimmt die Bindekorrektur
\setlength{\oddsidemargin}{2cm - 1in + 0.5cm}
\setlength{\textwidth}{\paperwidth - (1in + \hoffset) - \oddsidemargin - 4cm}
\setlength{\evensidemargin}{\paperwidth - (1in + \hoffset)*2 - \oddsidemargin - \textwidth}
\setlength{\marginparwidth}{4cm - \marginparsep{} - 1cm}

% Wenn größerer Zeilenabstand gewünscht ist, je nach Schriftgröße:
% Abstand        10pt    11pt    12pt
% -----------------------------------
% anderthalb     1.25    1.21    1.24
% doppelt        1.67    1.62    1.66
%
%\renewcommand{\baselinestretch}{1.21}

% Kopf- und Fußzeilen einrichten
\pagestyle{fancyplain}
\setlength{\headwidth}{\textwidth}
\addtolength{\headwidth}{\marginparsep}
\addtolength{\headwidth}{\marginparwidth}
\renewcommand{\chaptermark}[1]{\markboth{#1}{}}
\renewcommand{\sectionmark}[1]{\markright{\thesection\ #1}}
\lhead[\fancyplain{}{\bfseries\thepage}]%
	{\fancyplain{}{\bfseries\rightmark}}
\rhead[\fancyplain{}{\bfseries\leftmark}]%
	{\fancyplain{}{\bfseries\thepage}}
\cfoot{}

\begin{document}
% ---------------------------------- Vorspann ----------------------------------
% In den folgenden Zeilen werden die wichtigsten Informationen zur Arbeit ge-
% setzt, die dann im Dokument an den entsprechenden Stellen eingefügt werden.
%
\title{Titel der Arbeit}
\author{Name des Autoren}
\supervisor{Name des Referenten}
\cosupervisor{Name des Korreferenten}
\date{Datum der Abgabe}
\city{Ort}

% Folgende Befehle steuern, welchen Verbreitungsformen zugestimmt wird. Nicht
% gewünschte Verbreitungsformen entsprechend auskommentieren.
%
\publishlibrary  % Einstellung der Arbeit in die Bibliothek der FHW
\publishtitle    % Veröffentlichung des Titels der Arbeit im Internet
\publishdocument % Veröffentlichung der Arbeit im Internet

% Ab hier sollten keine Änderungen des Vorspanns nötig sein!
% -> Weiter zum Hauptteil

\pagenumbering{roman} % Bis zum ersten Kapitel mit römischen Seitenzahlen

\begin{titlepage}
	\begin{center}
		% Kopf der Seite
		\includegraphics{hsrm-logo} \\[0.7cm]
		{Fachbereich Design Informatik Medien}
		
		\vfill

		{\LARGE Abschlussarbeit} \\[0.5cm]
		{\large zur Erlangung des akademischen Grades} \\[0.5cm]
		{\large Bachelor of Science (B.Sc.)}
		
		% Titel
		\rule{\textwidth}{1pt}\\[0.5cm]
		{\huge \bfseries \thetitle}\\[0.1cm]
		\rule{\textwidth}{1pt}
		
		\vfill
		
		% Mitte der Seite
		\begin{tabular}{lr}
			Vorgelegt von & \theauthor \\
			am & \thedate \\
			Referent & \thesupervisor \\
			Korreferent & \thecosupervisor
		\end{tabular}
		
		\vfill
		
		% Fuß der Seite
		
	\end{center}
\end{titlepage}
 % Titelseite (wird automatisch mit Werten von oben gefüllt)
\clearemptydoublepage % Sorgt dafür, dass die nächste Seite rechts beginnt

\section*{Erklärung gemäß PO Teil A § 6.4.2}

Ich erkläre hiermit,
\begin{itemize}
	\item dass ich die vorliegende Abschlussarbeit selbstständig angefertigt,
	\item keine anderen als die angegebenen Quellen benutzt,
	\item die wörtlich oder dem Inhalt nach aus fremden Arbeiten entnommenen Stellen, bildlichen Darstellungen und dergleichen als solche genau kenntlich gemacht und
	\item keine unerlaubte fremde Hilfe in Anspruch genommen habe.
\end{itemize}

\vspace{3cm}
\makebox[.45\linewidth][l]{\thecity, \thedate}\hfill\rule{.45\linewidth}{0.5pt}\\
\makebox[.45\linewidth][l]{}\hfill\makebox[.45\linewidth][l]{\theauthor}

\vfill

\section*{Erklärung zur Verwendung der Diplomarbeit}

Hiermit erkläre ich mein Einverständnis mit den im folgenden aufgeführten Verbreitungsformen dieser Abschlussarbeit:

\newcolumntype{R}{>{\raggedright\arraybackslash}X}
\begin{tabularx}{\linewidth}{|R|c|c|}
	\hline
	\textbf{Verbreitungsform} & \textbf{Ja} & \textbf{Nein}
	\\\hline
	Einstellung der Arbeit in die Bibliothek der Hochschule RheinMain & \ifthenelse{\boolean{publishlibrary}}{$\times$}{} & \ifthenelse{\boolean{publishlibrary}}{}{$\times$}
	\\\hline
	Veröffentlichung des Titels der Arbeit im Internet & \ifthenelse{\boolean{publishtitle}}{$\times$}{} & \ifthenelse{\boolean{publishtitle}}{}{$\times$}
	\\\hline
	Veröffentlichung der Arbeit im Internet & \ifthenelse{\boolean{publishdocument}}{$\times$}{} & \ifthenelse{\boolean{publishdocument}}{}{$\times$}
	\\\hline
\end{tabularx}

\vspace{3cm}
\makebox[.45\linewidth][l]{\thecity, \thedate}\hfill\rule{.45\linewidth}{0.5pt}\\
\makebox[.45\linewidth][l]{}\hfill\makebox[.45\linewidth][l]{\theauthor}
 % Erklärung zur Prüfungsordnung und Verwendung der Arbeit
\clearemptydoublepage

\tableofcontents % Inhaltsverzeichnis generieren

\clearemptydoublepage
\pagenumbering{arabic}

% --------------------------------- Hauptteil ----------------------------------
% Im diesen Teil werden die einzelnen Kapitel eingefügt, die sinnvollerweise
% im Verzeichnis 'kapitel' abgelegt werden.

\chapter{Einleitung}

Hier Beginnt das Dokument \ldots

%\include{kapitel/zweileitung}

% und so weiter ...

% ---------------------------------- Anhänge -----------------------------------
% In diesem Teil werden alle Anhänge eingefügt, die auch als ganz normale Kapi-
% tel abgelegt werden.
\appendix

%\include{kapitel/einanhang}

% Ausgabe des Glossars (oder der Glossare, wenn mehrere definiert sind)
\printglossaries

% Ausgabe des Literaturverzeichnisses
\bibliographystyle{geralpha}
\bibliography{literatur}

\end{document}
